\newpage
\section{Introduction}

CineExplorer is an exploration platform that integrates the \ac{imdb} non-commercial dataset~\cite{imdb-noncommercial} with a knowledge-graph architecture. The aim is to move beyond the row-and-column queries that \ac{imdb} itself supports, towards graph-shaped questions that exploit the connections between persons, titles, and genres. The application leans on the standard Semantic Web stack: an \ac{owl}~2 QL ontology defines the conceptual model, an \ac{r2rml} mapping converts the relational source into \ac{rdf}, an Apache Fuseki triplestore serves \ac{sparql}~1.1 queries, and an Apache Brwsr instance dereferences resource \acp{iri} as Linked Data~\cite{cooluris}. Wikidata~\cite{wikidata} is integrated through federated \ac{sparql} \texttt{SERVICE} calls to enrich local persons with external attributes.

The report covers four milestones. Milestone~1 establishes the relational source: a 14-relation schema reduced from a conceptual \ac{erd}, populated from a numVotes-ranked top-$N$ slice of the official \ac{imdb} dataset (Sections~\ref{sec:dataset}--\ref{sec:infosystem}). Milestone~2 designs the ontology --- centred on \texttt{CreativeWork} as the abstraction of all titles, and a reified \texttt{Participation} class that lets per-credit context (role, job, character names) hang off the relation rather than the endpoints (Section~\ref{sec:ontology}). Milestone~3 connects the two layers with an \ac{r2rml} mapping that produces a 1{,}611{,}676-triple knowledge graph (Section~\ref{sec:mapping}). Milestone~4 deploys the result, evaluates ten \ac{sparql} queries that exercise \ac{sparql}~1.1's main constructs (subqueries, federated \texttt{SERVICE}, property paths, aggregates, negation as failure), and presents two non-trivial demonstrators built on top: a Bacon-number collaboration analysis and a \ac{shacl} structural validation (Sections~\ref{sec:deployment}--\ref{sec:demonstrator}).

Three application directions guided the design throughout. First, the system supports \emph{collaboration analysis}: graph-shaped questions of the form ``how are these two persons connected through shared work,'' formalised by the Bacon-number game and answered with \ac{sparql} property paths. Second, it supports \emph{enrichment via federation}: local persons are joined with Wikidata through their \ac{imdb} identifiers, without materialising any \texttt{owl:sameAs} bridge. Third, it supports \emph{Linked Data exploration}: every resource in the \ac{kg} is dereferenceable through Brwsr, so navigation can proceed without writing \ac{sparql}.
